\section{
    Tendências e Pesquisas em Andamento
} \label{sec:tendencias}


As tecnologias de desenvolvimento e as ameaças a aplicações web evoluem em ritmo semelhante, impulsionando pesquisas em novas abordagens de proteção. A mitigação de ataques complexos exige a incorporação contínua da segurança, do design à operação. A seguir, apresentam-se tendências e frentes de pesquisa atuais.

\subsection{
    Inteligência Artificial para Detecção de Ameaças
}

A inteligência artificial, especialmente o aprendizado de máquina, tem ampliado a capacidade de detecção de vulnerabilidades e de classificação de anomalias. Modelos como CNN, LSTM e GNN, combinados com técnicas como N-Grams e \textit{support vector machines} (SVM), aprimoram a identificação de padrões de injeção, incluindo XSS e SQL. Também se destacam abordagens comportamentais voltadas a ataques \textit{zero-day} e estudos de robustez adversarial. Como desafio central, permanece a interpretabilidade dos modelos em IDS, de modo que decisões de bloqueio sejam auditáveis por equipes de segurança.

\subsection{
    Arquiteturas de Confiança Zero e Gestão Avançada de Identidades
}

Com a consolidação de microsserviços e APIs distribuídas, arquiteturas de \textit{Zero Trust} tornam-se essenciais. Nesse modelo, nenhuma requisição é considerada confiável por padrão, exigindo verificação contínua de identidade, dispositivo e contexto, inclusive em tráfego interno. Como aplicações atravessam múltiplos domínios, faz-se necessária a federação de identidades e o estabelecimento dinâmico de confiança. Padrões como SAML 2.0 e XACML permanecem relevantes para SSO e autorização baseada em atributos, com uso de identificadores opacos para reduzir rastreabilidade e reforçar a privacidade.

\subsection{
    \textit{DevSecOps} - Integração de segurança no CI/CD
}

A adoção de \textit{DevSecOps} busca inserir segurança de forma nativa no ciclo de desenvolvimento, com verificações automatizadas de integridade e vulnerabilidades ao longo do CI/CD. A integração de SAST, DAST e IAST antecipa a identificação de falhas de injeção e erros lógicos antes da produção, reduzindo o custo de correção. Essa estratégia responde à necessidade de entregas ágeis sem perda de controle e reforça o princípio \textit{shift-left}, em consonância com o destaque de "design inseguro" na OWASP \textit{Top} 10 (2021).

\subsection{
    Segurança em Cadeias de Suprimento e Componentes de Terceiros
}

O uso intensivo de bibliotecas de terceiros, repositórios abertos e CDNs amplia os riscos na cadeia de suprimento de software. Pesquisas atuais priorizam o monitoramento contínuo de dependências, a validação de procedência e a verificação de integridade por assinaturas digitais e ferramentas como OWASP \textit{Dependency-Check}. O objetivo é impedir que atualizações automáticas ou pacotes comprometidos introduzam código malicioso e vulnerabilidades legadas no pipeline de CI/CD, fortalecendo a rastreabilidade e a resposta a incidentes.

\subsection{
    Desafios Abertos
}

Apesar dos avanços, permanecem desafios técnicos relevantes. O predomínio de APIs (REST e GraphQL), SPAs e processamento \textit{DOM-based} altera a superfície de ataque e dificulta a mitigação de falhas no cliente, como DOM-\textit{based} XSS, exigindo sanitização contextual mais robusta no navegador. Em paralelo, APIs ocultas (\textit{shadow} APIs) criam pontos não mapeados de risco e demandam descoberta contínua, políticas granulares de autorização e controles de taxa, uma vez que proteções estritamente no cliente podem ser contornadas por ataques diretos à API.
